% D5 boundary reset theorem: proof attempt 055

\paragraph{Boundary reset theorem (reduced uniform target).}
Let
\[
B=(s,u,v,\lambda,f),\qquad \tau(x)=\max\{t\ge 0: dn(F^j x)=\delta_{\mathrm{flat}}\text{ for }0\le j<t\},
\]
with \(\delta_{\mathrm{flat}}=(0,0,0,1)\). Then away from the boundary,
\[
\tau(Fx)=\tau(x)-1\qquad (\tau(x)>0).
\]
So to prove the uniform odd-\(m\) boundary theorem it is enough to prove the following branch lemmas at \(\tau(x)=0\):
\begin{enumerate}
\item \emph{Wrap branch:}
\[
\epsilon_4(x)=\mathrm{wrap}\implies \tau(Fx)=0.
\]
\item \emph{Carry\_jump branch:} there is a current-state formula
\[
\tau(Fx)=R_{\mathrm{cj}}(s,v,\lambda),
\]
with checked candidate
\[
R_{\mathrm{cj}}(s,v,\lambda)=
\begin{cases}
0,& s+v+\lambda\equiv 2\pmod m,\\
1,& s=1 \text{ and } s+v+\lambda\not\equiv 2\pmod m,\\
m-2,& \text{otherwise}.
\end{cases}
\]
Equivalently, on the checked range,
\[
q\equiv 1-s-v-\lambda\pmod m
\]
on the carry\_jump boundary.
\item \emph{Other branch:} there is a current-state formula
\[
\tau(Fx)=R_{\mathrm{oth}}(s,u,\lambda).
\]
On the checked range the branch splits into the two concrete grouped-delta subtypes
\[
(0,0,1,0)\qquad\text{and}\qquad(1,0,0,0),
\]
with reset values
\[
(0,0,1,0)\mapsto m-4,
\qquad
(1,0,0,0)\mapsto
\begin{cases}
m-3,& s=1,\\
0,& s\neq 1.
\end{cases}
\]
\end{enumerate}
Thus the full boundary theorem reduces to proving the branch lemmas uniformly in odd \(m\).
